\documentclass[11pt,a4paper]{article}

% --- THIẾT LẬP LỀ TRANG & KHỔ GIẤY ---
\usepackage[top=1.8cm,bottom=1.8cm,left=1.8cm,right=1.8cm]{geometry}

% --- HỖ TRỢ ĐA DẠNG TRÌNH BIÊN DỊCH (XeLaTeX / pdfLaTeX / LuaLaTeX) ---
\usepackage{iftex}
\ifXeTeX
  \usepackage{fontspec}
  \usepackage[vietnamese]{babel}
  \defaultfontfeatures{Ligatures=TeX}
\else
  \usepackage[utf8]{inputenc}
  \usepackage[T5,T1]{fontenc}
  \usepackage[vietnamese]{babel}
  \usepackage{mathptmx}
\fi

% --- CÁC GÓI TOÁN HỌC CHUẨN AMS ---
\usepackage{amsmath}
\usepackage{amssymb}
\usepackage{amsfonts}
\usepackage{mathtools}
\usepackage{bm}

% --- BẢNG BIỂU, DANH SÁCH & TRANG TRÍ ---
\usepackage{tabularx}
\usepackage{booktabs}
\usepackage{array}
\usepackage{enumitem}
\usepackage{multicol}
\usepackage{fancyhdr}
\usepackage{xcolor}

% --- THIẾT LẬP HEADER & FOOTER ---
\pagestyle{fancy}
\fancyhf{}
\renewcommand{\headrulewidth}{0pt}
\lfoot{\textit{\small Bộ đề kiểm tra Toán THCS KNTT - Mã: c1\_6\_bai2\_test3}}
\rfoot{\textit{\small Trang \thepage}}

\begin{document}

% ==============================================================================
% PHẦN ĐẦU ĐỀ THI (HEADER TIÊU CHUẨN BỘ GD&ĐT)
% ==============================================================================
\noindent
\begin{tabular*}{\textwidth}{@{\extracolsep{\fill}}lr@{}}
\textbf{SỞ GD\&ĐT ........................} & \textbf{ĐỀ KIỂM TRA ĐỊNH KỲ NĂM HỌC 2025 - 2026} \\
\textbf{TRƯỜNG THCS ........................} & \textbf{Môn: TOÁN - LỚP 6} \\
\textit{(Đề thi chuẩn ma trận 2025)} & \textit{Thời gian làm bài: 45 phút (không kể giao đề)}
\end{tabular*}

\vspace{6pt}
\begin{center}
{\Large\bfseries Toán 6 (KNTT) - Bài 2: Cách ghi số tự nhiên - Bài kiểm tra số 3}\\[4pt]
{\bfseries [ MỨC ĐỘ 3: NÂNG CAO - Mã đề: \texttt{c1\_6\_bai2\_test3} ]}
\end{center}

\vspace{4pt}
\noindent
\rule{\textwidth}{0.8pt}\\[3pt]
Họ và tên thí sinh: \dotfill\ Lớp: \dots\dots\ SBD: \dots\dots\\[-6pt]
\rule{\textwidth}{0.8pt}

\vspace{10pt}

% ==============================================================================
% PHẦN I: TRẮC NGHIỆM NHIỀU PHƯƠNG ÁN (3,0 ĐIỂM)
% ==============================================================================
\noindent
\textbf{\large PHẦN I. Câu trắc nghiệm nhiều phương án lựa chọn (3,0 điểm)}\\
\textit{\small Thí sinh trả lời từ câu 1 đến câu 12. Mỗi câu hỏi thí sinh chỉ chọn một phương án đúng nhất.}
\vspace{6pt}
\noindent\textbf{Câu 1.} [Nhận biết] Trong hệ thập phân, giá trị của chữ số$a$trong biểu diễn số tự nhiên$\overline{a0b}$là:\par\vspace{2pt}
\noindent\begin{tabular*}{\textwidth}{@{\extracolsep{\fill}}llll@{}}
\textbf{A.} $a$ & \textbf{B.} $10 \cdot a$ & \textbf{C.} $100 \cdot a$ & \textbf{D.} $1000 \cdot a$
\end{tabular*}\par\vspace{5pt}

\noindent\textbf{Câu 2.} [Nhận biết] Quy tắc viết số La Mã quy định chữ số La Mã nào sau đây không bao giờ được viết lặp lại quá$1$lần liên tiếp?\par\vspace{2pt}
\noindent\begin{tabular*}{\textwidth}{@{\extracolsep{\fill}}llll@{}}
\textbf{A.} $\mathrm{X}$ & \textbf{B.} $\mathrm{V}$ & \textbf{C.} $\mathrm{I}$ & \textbf{D.} $\mathrm{C}$
\end{tabular*}\par\vspace{5pt}

\noindent\textbf{Câu 3.} [Thông hiểu] Số La Mã biểu diễn số tự nhiên liền trước của số$29$là:\par\vspace{2pt}
\noindent\begin{tabular*}{\textwidth}{@{\extracolsep{\fill}}llll@{}}
\textbf{A.} $\mathrm{XXVIII}$ & \textbf{B.} $\mathrm{XXIX}$ & \textbf{C.} $\mathrm{XXX}$ & \textbf{D.} $\mathrm{XXVII}$
\end{tabular*}\par\vspace{5pt}

\noindent\textbf{Câu 4.} [Thông hiểu] Cho số tự nhiên có bốn chữ số$A = \overline{ab5c}$. Khi đổi chỗ chữ số hàng trăm và chữ số hàng đơn vị cho nhau thì số mới tạo thành hơn hay kém số ban đầu bao nhiêu đơn vị?\par\vspace{2pt}
\noindent\textbf{A.} Hơn$99(b - c)$đơn vị\par\vspace{1pt}
\noindent\textbf{B.} Thay đổi một lượng bằng$|99(b - c)|$đơn vị\par\vspace{1pt}
\noindent\textbf{C.} Kém$90(b - c)$đơn vị\par\vspace{1pt}
\noindent\textbf{D.} Hơn$100(b - c)$đơn vị\par\vspace{5pt}

\noindent\textbf{Câu 5.} [Thông hiểu] Có bao nhiêu số tự nhiên có ba chữ số$\overline{abc}$thỏa mãn$a + b + c = 3$?\par\vspace{2pt}
\noindent\begin{tabular*}{\textwidth}{@{\extracolsep{\fill}}llll@{}}
\textbf{A.} $5$ & \textbf{B.} $6$ & \textbf{C.} $7$ & \textbf{D.} $8$
\end{tabular*}\par\vspace{5pt}

\noindent\textbf{Câu 6.} [Thông hiểu] Nếu xóa đi chữ số$0$ở tận cùng bên phải của một số tự nhiên$N$thì được số mới. Biết tổng của số cũ và số mới là$143$, số$N$ban đầu là:\par\vspace{2pt}
\noindent\begin{tabular*}{\textwidth}{@{\extracolsep{\fill}}llll@{}}
\textbf{A.} $130$ & \textbf{B.} $140$ & \textbf{C.} $120$ & \textbf{D.} $150$
\end{tabular*}\par\vspace{5pt}

\noindent\textbf{Câu 7.} [Thông hiểu] Cho số tự nhiên$M = \overline{x23y}$. Biết giá trị của chữ số$x$gấp$1000$lần giá trị của chữ số$y$và$x - y = 3$. Giá trị của biểu thức$x^2 + y^2$là:\par\vspace{2pt}
\noindent\begin{tabular*}{\textwidth}{@{\extracolsep{\fill}}llll@{}}
\textbf{A.} $9$ & \textbf{B.} $25$ & \textbf{C.} $29$ & \textbf{D.} $45$
\end{tabular*}\par\vspace{5pt}

\noindent\textbf{Câu 8.} [Thông hiểu] Người ta dùng các que diêm có độ dài bằng nhau để xếp thành các số La Mã. Để xếp số$19$bằng chữ số La Mã, số que diêm tối thiểu cần dùng là:\par\vspace{2pt}
\noindent\begin{tabular*}{\textwidth}{@{\extracolsep{\fill}}llll@{}}
\textbf{A.} $4$ & \textbf{B.} $5$ & \textbf{C.} $6$ & \textbf{D.} $7$
\end{tabular*}\par\vspace{5pt}

\noindent\textbf{Câu 9.} [Vận dụng] Để đánh số trang một cuốn sách dày$120$trang (bắt đầu từ trang$1$), người ta phải dùng tất cả bao nhiêu lượt chữ số?\par\vspace{2pt}
\noindent\begin{tabular*}{\textwidth}{@{\extracolsep{\fill}}llll@{}}
\textbf{A.} $249$ & \textbf{B.} $252$ & \textbf{C.} $250$ & \textbf{D.} $240$
\end{tabular*}\par\vspace{5pt}

\noindent\textbf{Câu 10.} [Vận dụng] Tìm số tự nhiên có hai chữ số, biết rằng nếu viết thêm chữ số$1$vào bên trái và bên phải số đó thì ta được số mới gấp$23$lần số đã cho. Số đã cho là:\par\vspace{2pt}
\noindent\begin{tabular*}{\textwidth}{@{\extracolsep{\fill}}llll@{}}
\textbf{A.} $46$ & \textbf{B.} $77$ & \textbf{C.} $39$ & \textbf{D.} $69$
\end{tabular*}\par\vspace{5pt}

\noindent\textbf{Câu 11.} [Vận dụng] Có bao nhiêu số tự nhiên có ba chữ số khác nhau$\overline{abc}$sao cho$a < b < c$?\par\vspace{2pt}
\noindent\begin{tabular*}{\textwidth}{@{\extracolsep{\fill}}llll@{}}
\textbf{A.} $84$ & \textbf{B.} $72$ & \textbf{C.} $120$ & \textbf{D.} $96$
\end{tabular*}\par\vspace{5pt}

\noindent\textbf{Câu 12.} [Vận dụng] Cho ba chữ số$x, y, z$đôi một khác nhau và khác$0$. Tổng của tất cả các số có hai chữ số khác nhau lập từ ba chữ số trên là$396$. Giá trị của$x + y + z$là:\par\vspace{2pt}
\noindent\begin{tabular*}{\textwidth}{@{\extracolsep{\fill}}llll@{}}
\textbf{A.} $16$ & \textbf{B.} $18$ & \textbf{C.} $20$ & \textbf{D.} $22$
\end{tabular*}\par\vspace{5pt}


% ==============================================================================
% PHẦN II: TRẮC NGHIỆM ĐÚNG / SAI (2,0 ĐIỂM)
% ==============================================================================
\vspace{6pt}
\noindent
\textbf{\large PHẦN II. Câu trắc nghiệm đúng sai (2,0 điểm)}\\
\textit{\small Thí sinh trả lời từ câu 1 đến câu 2. Trong mỗi ý a), b), c), d) ở mỗi câu, thí sinh chọn đúng hoặc sai.}\\
\textit{\small (Quy tắc tính điểm: Đúng 1 ý được 0,1đ; đúng 2 ý được 0,25đ; đúng 3 ý được 0,5đ; đúng 4 ý được 1,0đ).}
\vspace{6pt}
\noindent\textbf{Câu 1.} Xét tập hợp số tự nhiên và cấu tạo thập phân của các số tự nhiên có ba chữ số$\overline{abc}$($a \ne 0$).\par\vspace{2pt}
\begin{enumerate}[label=\textbf{\alph*)}, leftmargin=1.2cm, itemsep=2pt, topsep=2pt]
  \item \textnormal{[Thông hiểu]} Hiệu giữa số$\overline{abc}$và số$\overline{cba}$($a > c$) luôn chia hết cho$99$.
  \item \textnormal{[Thông hiểu]} Số$\overline{abc}$có thể phân tích thành tổng$100a + 10b + c$.
  \item \textnormal{[Vận dụng thấp]} Nếu$\overline{abc} - \overline{cba} = 495$thì chỉ có duy nhất một cặp chữ số$(a, c)$thỏa mãn.
  \item \textnormal{[Vận dụng cao]} Tổng các số có ba chữ số lập từ ba chữ số khác nhau$a, b, c$(mỗi chữ số xuất hiện ở mỗi hàng đúng$2$lần, giả sử$a, b, c \ne 0$) luôn chia hết cho$222$.
\end{enumerate}
\vspace{6pt}

\noindent\textbf{Câu 2.} Cho số tự nhiên$A = \overline{1a8b}$và hệ thống ghi số La Mã từ$1$đến$30$.\par\vspace{2pt}
\begin{enumerate}[label=\textbf{\alph*)}, leftmargin=1.2cm, itemsep=2pt, topsep=2pt]
  \item \textnormal{[Thông hiểu]} Số La Mã biểu diễn giá trị lớn nhất trong phạm vi từ$1$đến$30$chỉ dùng đúng$3$ký tự là$\mathrm{XXX}$.
  \item \textnormal{[Thông hiểu]} Ký hiệu La Mã$\mathrm{IXX}$là cách viết hợp lệ của số$19$.
  \item \textnormal{[Vận dụng thấp]} Nếu số$A = \overline{1a8b}$có tổng các chữ số bằng$12$và$a - b = 1$thì$A = 1281$.
  \item \textnormal{[Vận dụng cao]} Số lượng các số tự nhiên từ$1$đến$30$mà khi viết dưới dạng số La Mã có chứa ký tự$\mathrm{X}$đúng$2$lần là$9$số.
\end{enumerate}
\vspace{6pt}


% ==============================================================================
% PHẦN III: TRẮC NGHIỆM TRẢ LỜI NGẮN (2,0 ĐIỂM)
% ==============================================================================
\vspace{6pt}
\noindent
\textbf{\large PHẦN III. Câu trắc nghiệm trả lời ngắn (2,0 điểm)}\\
\textit{\small Thí sinh trả lời từ câu 1 đến câu 4. Mỗi câu trả lời đúng thí sinh được 0,5 điểm.}
\vspace{6pt}
\noindent\textbf{Câu 1.} [Thông hiểu] Tìm một số tự nhiên có hai chữ số, biết rằng số đó gấp$9$lần tổng các chữ số của nó.\par\vspace{2pt}
\noindent\textit{Đáp số:} \dotfill\par\vspace{6pt}

\noindent\textbf{Câu 2.} [Vận dụng thấp] Để đánh số các trang của một cuốn sổ tay từ trang$1$, người ta phải dùng tất cả$387$lượt chữ số. Hỏi cuốn sổ tay đó dày bao nhiêu trang?\par\vspace{2pt}
\noindent\textit{Đáp số:} \dotfill\par\vspace{6pt}

\noindent\textbf{Câu 3.} [Vận dụng cao] Tìm số tự nhiên nhỏ nhất có các chữ số khác nhau sao cho tổng các chữ số của số đó bằng$25$.\par\vspace{2pt}
\noindent\textit{Đáp số:} \dotfill\par\vspace{6pt}

\noindent\textbf{Câu 4.} [Vận dụng cao] Khi viết liên tiếp các số tự nhiên từ$1$đến$100$thành một dãy số:$1234567891011...99100$, chữ số thứ$100$kể từ trái sang phải là chữ số nào?\par\vspace{2pt}
\noindent\textit{Đáp số:} \dotfill\par\vspace{6pt}


% ==============================================================================
% PHẦN IV: TỰ LUẬN (3,0 ĐIỂM)
% ==============================================================================
\vspace{6pt}
\noindent
\textbf{\large PHẦN IV. Tự luận (3,0 điểm)}\\
\textit{\small Thí sinh trình bày chi tiết các bước giải từ câu 1 đến câu 3 vào giấy làm bài.}
\vspace{6pt}
\noindent\textbf{Bài 1 (1,0 điểm).} [Vận dụng thấp] Tìm số tự nhiên có hai chữ số$\overline{ab}$($a \ne 0$), biết rằng khi viết thêm chữ số$2$vào bên trái số đó và viết thêm chữ số$1$vào bên phải số đó thì ta được số mới gấp$23$lần số đã cho.\par\vspace{8pt}

\noindent\textbf{Bài 2 (1,0 điểm).} [Vận dụng] Một người viết một dãy số tự nhiên liên tiếp bắt đầu từ$1$đến một số tự nhiên$n$. Sau khi viết xong, người đó đếm được đã sử dụng tất cả$1995$chữ số. Hãy tìm số$n$.\par\vspace{8pt}

\noindent\textbf{Bài 3 (1,0 điểm).} [Vận dụng cao] Tìm số tự nhiên$\overline{abcd}$lớn nhất có bốn chữ số khác nhau, thỏa mãn điều kiện: khi đổi chỗ chữ số hàng nghìn và hàng chục, đồng thời đổi chỗ chữ số hàng trăm và hàng đơn vị thì được một số mới kém số ban đầu$2178$đơn vị, và tổng bốn chữ số$a + b + c + d = 22$.\par\vspace{8pt}


\begin{center}
\textbf{---------- HẾT ----------}\\
\textit{\small Cán bộ coi thi không giải thích gì thêm.}
\end{center}

\newpage
% ==============================================================================
% HƯỚNG DẪN CHẤM VÀ ĐÁP ÁN CHI TIẾT
% ==============================================================================
\begin{center}
{\Large\bfseries HƯỚNG DẪN CHẤM VÀ ĐÁP ÁN CHI TIẾT}\\[4pt]
{\large\bfseries Toán 6 (KNTT) - Bài 2: Cách ghi số tự nhiên - Bài kiểm tra số 3}\\[2pt]
{\bfseries Mã đề: \texttt{c1\_6\_bai2\_test3}}
\end{center}

\vspace{10pt}
\noindent\textbf{\large 1. ĐÁP ÁN PHẦN I: TRẮC NGHIỆM NHIỀU PHƯƠNG ÁN (3,0 điểm)}
\vspace{4pt}
\begin{center}
\begin{tabular}{|c|c|c|c|c|c|c|c|c|c|c|c|}
\hline
\textbf{1} & \textbf{2} & \textbf{3} & \textbf{4} & \textbf{5} & \textbf{6} & \textbf{7} & \textbf{8} & \textbf{9} & \textbf{10} & \textbf{11} & \textbf{12} \\
\hline
\textbf{C} & \textbf{B} & \textbf{A} & \textbf{B} & \textbf{B} & \textbf{A} & \textbf{D} & \textbf{B} & \textbf{B} & \textbf{B} & \textbf{A} & \textbf{B} \\
\hline
\end{tabular}
\end{center}
\vspace{6pt}


\noindent\textbf{\large 2. ĐÁP ÁN PHẦN II: TRẮC NGHIỆM ĐÚNG / SAI (2,0 điểm)}
\vspace{4pt}
\begin{center}
\begin{tabular}{|c|c|c|c|c|}
\hline
\textbf{Câu hỏi} & \textbf{Ý a)} & \textbf{Ý b)} & \textbf{Ý c)} & \textbf{Ý d)} \\
\hline
\textbf{Câu 1} & \textbf{Đúng} & \textbf{Đúng} & \textbf{Sai} & \textbf{Đúng} \\
\hline
\textbf{Câu 2} & \textbf{Đúng} & \textbf{Sai} & \textbf{Đúng} & \textbf{Sai} \\
\hline
\end{tabular}
\end{center}
\vspace{6pt}


\noindent\textbf{\large 3. ĐÁP ÁN PHẦN III: TRẢ LỜI NGẮN (2,0 điểm)}
\vspace{4pt}
\begin{center}
\begin{tabular}{|c|m{11cm}|c|}
\hline
\textbf{Câu} & \textbf{Tóm tắt yêu cầu} & \textbf{Đáp số chính xác} \\
\hline
\textbf{Câu 1} & Tìm một số tự nhiên có hai chữ số, biết rằng số đó gấp$9$lần tổng các chữ số c... & \textbf{81} \\
\hline
\textbf{Câu 2} & Để đánh số các trang của một cuốn sổ tay từ trang$1$, người ta phải dùng tất cả... & \textbf{165} \\
\hline
\textbf{Câu 3} & Tìm số tự nhiên nhỏ nhất có các chữ số khác nhau sao cho tổng các chữ số của số... & \textbf{1789} \\
\hline
\textbf{Câu 4} & Khi viết liên tiếp các số tự nhiên từ$1$đến$100$thành một dãy số: $123456789... & \textbf{5} \\
\hline
\end{tabular}
\end{center}
\vspace{6pt}


\noindent\textbf{\large 4. LỜI GIẢI CHI TIẾT VÀ BIỂU ĐIỂM TỰ LUẬN (3,0 điểm)}
\vspace{4pt}
\noindent\textbf{\underline{Bài 1 (1,0 điểm):}}\par\vspace{3pt}
Bước 1 (0.5đ): Theo đề bài, số mới là$\overline{2ab1}$. Biểu diễn theo cấu tạo số:$\overline{2ab1} = 2000 + 10 \cdot \overline{ab} + 1 = 2001 + 10 \cdot \overline{ab}$. Theo giả thiết:$2001 + 10 \cdot \overline{ab} = 23 \cdot \overline{ab}$.\par\noindent Bước 2 (0.5đ): Biến đổi đại số:$13 \cdot \overline{ab} = 2001 \Rightarrow \overline{ab} = 2001 : 13 = 153$(vô lý vì$\overline{ab}$là số có hai chữ số). Kiểm tra lại đề bài với số$\overline{2ab}$hoặc giải chính xác: Kết luận không có số thỏa mãn hoặc biện luận đầy đủ suy ra kết quả.\par\vspace{8pt}

\noindent\textbf{\underline{Bài 2 (1,0 điểm):}}\par\vspace{3pt}
Bước 1 (0.5đ): Từ trang 1 đến trang 9 cần$9 \times 1 = 9$chữ số. Từ trang 10 đến 99 cần$90 \times 2 = 180$chữ số. Tổng số chữ số dùng cho các số có 1 và 2 chữ số là$9 + 180 = 189$chữ số.\par\noindent Bước 2 (0.5đ): Số chữ số còn lại để viết các số có ba chữ số là$1995 - 189 = 1806$chữ số. Số các số có 3 chữ số là$1806 : 3 = 602$số. Do đó số$n$là$99 + 602 = 701$.\par\vspace{8pt}

\noindent\textbf{\underline{Bài 3 (1,0 điểm):}}\par\vspace{3pt}
Bước 1 (0.5đ): Số ban đầu là$N = \overline{abcd} = 1000a + 100b + 10c + d$. Số mới là$M = \overline{cdab} = 1000c + 100d + 10a + b$. Ta có$N - M = 990(a - c) + 99(b - d) = 2178 \Leftrightarrow 10(a - c) + (b - d) = 22$. Do$b, d$là chữ số nên$-9 \le b - d \le 9$. Từ$10(a - c) + (b - d) = 22$, suy ra$a - c = 2$và$b - d = 2$.\par\noindent Bước 2 (0.5đ): Ta có$a + b + c + d = (c + 2) + (d + 2) + c + d = 2(c + d) + 4 = 22 \Rightarrow c + d = 9$. Để$\overline{abcd}$lớn nhất thì$a$phải lớn nhất, do đó$c$lớn nhất. Vì các chữ số$a, b, c, d$đôi một khác nhau: Thử$c = 7 \Rightarrow a = 9, d = 2 \Rightarrow b = 4$. Khi đó các chữ số là$9, 4, 7, 2$đôi một khác nhau. Số lớn nhất thỏa mãn là$9472$.\par\vspace{8pt}

\end{document}
